\newpage%\新一页
\section{Assumptions and Justifications}
In response to the title of this article, the following hypotheses are proposed:
%1.	假设题目所给的数据真实可靠；
%注意：假设对整篇文章具有指导性，有时决定问题的难易。一定要注意假设的某种角度上的合理性，不能乱编，完全偏离事实或与题目要求相抵触。注意罗列要工整。

\begin{enumerate}[
	label = (\arabic*),
	itemindent = 0pt,
	labelindent = \parindent,
	labelwidth = 2em,
	labelsep = 5pt,
	leftmargin = 8em]
	\item[\textbf{Assumption 1:}] This is the first assumption 1.
	\item[\textbf{Assumption 2:}] This is the second assumption 2.
	\item[\textbf{Assumption 3:}] 
	\item[\textbf{Assumption 4:}] 
	\item[\textbf{Assumption 5:}] 
\end{enumerate}


\begin{enumerate}
	\item 
\end{enumerate}



\section{Notations}
%这部分不要过页（删掉此句话）
\begin{center}%置中
	\begin{tabular}{cc}
		\toprule[1pt] 
		\makebox[0.15\textwidth][c]{符号} & \makebox[0.4\textwidth][c]{说明} \\  
		\hline
		$T_i$&小温区温度\\
		$T_i$&小温区温度\\    
		\bottomrule[1pt]
		%尽可能借鉴参考书上通常采用的符号，不宜自己乱定义符号，对于改进的一些模型，符号可以适当自己修正（下标、上标、参数等可以变，主符号最好与经典模型符号靠近）。
		%对文章自己创新的名词需要特别解释。其他符号要进行说明，注意罗列要工整。如“ ～第 种疗法的第 项指标值”等，注意格式统一，不要出现零乱或前后不一致现象，关键是容易看懂。建议采用表格形式说明。
		注：未申明的变量以其在符号出现处的具体说明为准。
	\end{tabular}
\end{center}


